\rtmtight
\begin{tblr}{width=\textwidth,colspec={|Q[c,m,2.1cm]|X[l,m]|},hlines,vlines,hspan=minimal,rowsep=1.5pt,colsep=3pt,row{1}={bg=headgray,font=\bfseries}}
\SetCell{c,m}\hfil\logo\hfil & \textbf{POLITEKNIK NEGERI MALANG}\par\textbf{JURUSAN TEKNOLOGI INFORMASI}\par\textbf{PROGRAM STUDI: D4 TEKNIK INFORMATIKA} \\
\SetCell[c=2]{l} \textbf{RENCANA TUGAS MAHASISWA (RTM) DAN RUBRIK PENILAIAN} & \\
Nama Mata Kuliah & Metode Numerik \\
Kode Mata Kuliah & RTI253006 \quad \textbf{sks / Jam perminggu:} 2 SKS/4 jam \quad \textbf{Semester:} 3 \\
Dosen Pengampu & \begin{enumerate}\item Adevian Fairuz Pratama, S.ST., M.Eng\item Triana Fatmawati, S.T., M.T.\item Dian Hanifudin Subhi, S.Kom., M.Kom.\item M. Hasyim Ratsanjani, S.Kom., M.Kom.\end{enumerate} \\
\end{tblr}
\assessmentblock{Tugas Penyelesaian Masalah Numerik}{Tugas terstruktur individu}{SCPMK0707-03501, SCPMK0902-03502, SCPMK0902-03503, SCPMK1007-03504}{Mahasiswa menyelesaikan tugas perhitungan numerik secara individu mencakup analisis galat, penyelesaian sistem persamaan linear, persamaan nonlinear, diferensiasi, integrasi, interpolasi, dan regresi sesuai jadwal yang ditetapkan.}{Mempelajari konsep dan algoritma metode numerik, menyelesaikan soal perhitungan step-by-step secara manual dan/atau menggunakan perangkat lunak, mendokumentasikan proses dan hasil, kemudian mengumpulkan sesuai batas waktu.}{Lembar jawaban dengan proses perhitungan lengkap (kerja tangan dan/atau file perangkat lunak), tabel iterasi bila diperlukan, dan analisis galat.}{Ketepatan proses dan hasil perhitungan numerik & 8 & Proses perhitungan dilakukan dengan algoritma yang benar, tabel iterasi (bila ada) lengkap dan akurat, serta hasil akhir tepat sesuai kunci jawaban. & Proses perhitungan sebagian besar benar, tetapi ada kesalahan pada satu atau dua langkah iterasi atau hasil akhir kurang presisi. & Proses perhitungan salah secara fundamental, tabel iterasi tidak lengkap, atau hasil akhir sangat menyimpang dari nilai yang benar. \\ \\
Analisis galat dan konvergensi & 7 & Galat (absolut, relatif, pemotongan) dihitung dengan benar, kriteria konvergensi diterapkan tepat, dan analisis galat menunjukkan pemahaman yang baik. & Analisis galat sebagian besar benar, tetapi ada kesalahan pada jenis galat tertentu atau kriteria konvergensi belum diterapkan dengan tepat. & Analisis galat tidak dilakukan, salah secara konsep, atau tidak terhubung dengan proses iterasi. \\ \\
Kelengkapan dokumentasi dan keteraturan penyajian & 5 & Semua langkah ditulis dengan rapi dan teratur, notasi matematis benar, dan proses mudah diikuti oleh pembaca. & Langkah-langkah ditulis, tetapi ada beberapa notasi yang tidak konsisten atau penyajian yang kurang rapi. & Penyajian tidak teratur, banyak langkah yang terlewat, atau notasi matematis banyak yang salah. \\}

\assessmentblock{Kuis 1 Konsep Numerik dan Sistem Persamaan Linear}{Kuis (ujian tertulis, closed book)}{SCPMK0707-03501}{Mahasiswa mengerjakan Kuis 1 materi minggu 1--4 meliputi konsep metode numerik, analisis galat, dan penyelesaian sistem persamaan linear (Gauss, Gauss-Jordan, Gauss-Seidel) secara mandiri.}{Mengerjakan soal teori dan perhitungan numerik secara individu dengan sifat closed book sesuai jadwal asesmen.}{Lembar jawaban kuis dengan proses perhitungan lengkap.}{Penguasaan konsep metode numerik dan galat & 4 & Perbedaan metode analitik-numerik dijelaskan dengan benar, jenis galat diidentifikasi tepat, dan perhitungan galat absolut/relatif sesuai kunci. & Konsep dasar dipahami, tetapi ada kekeliruan pada jenis galat atau perhitungan galat kurang presisi. & Konsep metode numerik atau galat tidak dipahami, atau perhitungan sangat menyimpang dari kunci. \\ \\
Ketepatan penyelesaian sistem persamaan linear & 4 & Eliminasi Gauss, Gauss-Jordan, dan Gauss-Seidel diterapkan dengan benar; tabel iterasi lengkap dan solusi akurat. & Metode diterapkan sebagian besar benar, tetapi ada kesalahan pada pivoting, iterasi tertentu, atau solusi kurang presisi. & Metode tidak diterapkan dengan benar atau solusi sangat menyimpang dari nilai yang benar. \\ \\
Analisis dan pemilihan metode & 2 & Mahasiswa mampu membandingkan efisiensi metode dan memilih metode yang tepat sesuai karakteristik sistem persamaan yang diberikan. & Perbandingan metode dilakukan, tetapi argumen pemilihan belum didukung analisis yang cukup. & Tidak ada analisis perbandingan atau pemilihan metode tidak berdasar. \\}

\assessmentblock{UTS Metode Penyelesaian Persamaan}{UTS (ujian tertulis, closed book)}{SCPMK0707-03501, SCPMK0902-03502}{Mahasiswa mengerjakan soal UTS materi minggu 1--8 meliputi galat, sistem persamaan linear, dan persamaan nonlinear secara mandiri.}{Mengerjakan soal teori dan perhitungan numerik secara individu dengan sifat closed book sesuai jadwal asesmen.}{Lembar jawaban UTS dengan proses perhitungan lengkap.}{Penguasaan galat dan penyelesaian SPL & 10 & Analisis galat, eliminasi Gauss, Gauss-Jordan, dan Gauss-Seidel diterapkan benar dengan solusi yang akurat sesuai kunci. & Sebagian besar benar, tetapi ada kesalahan pada langkah tertentu atau galat analisis yang kurang tepat. & Banyak kesalahan fundamental pada galat atau penyelesaian SPL tidak dapat diselesaikan. \\ \\
Ketepatan metode pencarian akar persamaan nonlinear & 9 & Biseksi, regula falsi, Newton-Raphson, dan sekans diterapkan benar; tabel iterasi lengkap dan kriteria berhenti diterapkan tepat. & Metode diterapkan sebagian besar benar, tetapi ada kesalahan pada tabel iterasi atau kriteria berhenti. & Metode tidak dapat diterapkan dengan benar atau tabel iterasi sangat tidak lengkap. \\ \\
Analisis konvergensi dan pemilihan metode & 6 & Laju konvergensi metode dianalisis dengan benar, kondisi divergensi diidentifikasi, dan pemilihan metode didukung alasan teknis yang kuat. & Analisis konvergensi sebagian besar benar, tetapi pemilihan metode tidak selalu disertai alasan yang cukup. & Tidak ada analisis konvergensi atau pemilihan metode tidak berdasar. \\}

\assessmentblock{Kuis 2 Diferensiasi dan Integrasi Numerik}{Kuis (ujian tertulis, closed book)}{SCPMK0902-03503}{Mahasiswa mengerjakan Kuis 2 materi minggu 10--11 meliputi diferensiasi numerik dan integrasi Riemann secara mandiri.}{Mengerjakan soal teori dan perhitungan diferensiasi dan integrasi numerik secara individu dengan sifat closed book sesuai jadwal asesmen.}{Lembar jawaban kuis dengan proses perhitungan lengkap.}{Ketepatan perhitungan diferensiasi numerik & 4 & Diferensiasi maju, mundur, dan tengah dihitung benar; tingkat akurasi dan galat pemotongan dianalisis dengan tepat. & Perhitungan diferensiasi sebagian besar benar, tetapi analisis galat pemotongan belum lengkap. & Perhitungan diferensiasi salah secara fundamental atau galat tidak dianalisis. \\ \\
Ketepatan perhitungan integrasi Riemann & 4 & Integrasi Riemann kiri, kanan, dan tengah dihitung benar; perbandingan dengan nilai eksak dan analisis galat dilakukan dengan tepat. & Perhitungan Riemann sebagian besar benar, tetapi perbandingan atau analisis galat belum lengkap. & Metode Riemann tidak dapat diterapkan atau hasil sangat menyimpang dari nilai eksak. \\ \\
Kemampuan analisis akurasi & 2 & Mahasiswa mampu membandingkan akurasi diferensiasi maju, mundur, tengah dan Riemann kiri, kanan, tengah serta memberikan rekomendasi metode yang paling akurat. & Perbandingan dilakukan, tetapi rekomendasi belum didukung analisis galat yang memadai. & Tidak ada analisis perbandingan akurasi antar metode. \\}

\assessmentblock{UAS Implementasi Metode Numerik Lengkap}{UAS (ujian tertulis, closed book)}{SCPMK0902-03503, SCPMK1007-03504}{Mahasiswa mengerjakan soal UAS materi minggu 1--16 dengan penekanan pada integrasi numerik (trapezoidal, Simpson), interpolasi, dan regresi secara mandiri.}{Mengerjakan soal teori dan perhitungan numerik komprehensif secara individu dengan sifat closed book sesuai jadwal asesmen.}{Lembar jawaban UAS dengan proses perhitungan lengkap.}{Ketepatan integrasi numerik trapezoidal dan Simpson & 8 & Aturan trapezoidal komposit dan Simpson 1/3 serta 3/8 diterapkan benar; galat estimasi dihitung dan perbandingan akurasi antar metode dilakukan. & Metode diterapkan sebagian besar benar, tetapi estimasi galat atau perbandingan akurasi belum lengkap. & Banyak kesalahan pada penerapan trapezoidal atau Simpson atau tidak dapat mengestimasi galat. \\ \\
Ketepatan interpolasi dan analisis galat & 7 & Interpolasi Lagrange dan Newton dihitung benar; derajat polinomial dipilih tepat dan galat interpolasi dianalisis dengan benar. & Interpolasi sebagian besar benar, tetapi pilihan derajat polinomial atau analisis galat belum optimal. & Interpolasi tidak dapat diselesaikan atau galat tidak dianalisis. \\ \\
Ketepatan regresi dan interpretasi model & 5 & Persamaan regresi (linear dan nonlinear) dihitung benar menggunakan least squares, $R^2$ dihitung tepat, dan model diinterpretasikan secara bermakna. & Persamaan regresi sebagian besar benar, tetapi $R^2$ tidak dihitung atau interpretasi model kurang tepat. & Persamaan regresi salah secara fundamental atau tidak ada interpretasi model. \\}

\begin{tblr}{width=\textwidth,colspec={|Q[l,m,3.2cm]|X[l,m]|},hlines,vlines,hspan=minimal,row{1-3}={bg=headgray},rowsep=1.5pt,colsep=3pt}
Jadwal Pelaksanaan: & Minggu 1--17 sesuai RPS. Setiap evaluasi memiliki RTM dan rubrik multi-kriteria. \\
Lain-Lain Yang Diperlukan: & LMS, MS Excel/Python/MATLAB, kalkulator ilmiah, dan perangkat presentasi. \\
Pustaka: & Mengacu pustaka utama dan pendukung pada bagian RPS. \\
\end{tblr}
